\section{Implications and limits}
A school platform looking only at recorded offers would endorse adding rank to its predictor. A platform looking only at placement would reject the same change. Neither verdict mentions the increase in first choices or identifies the people who bear the loss. The families whose menus changed tended to move up in the replay. Families with the same menu who moved at all tended to move down. It is possible to value first choices highly and still want to protect those applicants. The calculation in Section~\ref{sec:row-screen} shows exactly what protecting every old contract permits. Appendix~\ref{sec:remedy} considers a narrower intervention after offers: referring an unmatched person to a vacancy without taking anyone else's contract.

The rental experiment addresses a different social concern. A renter who cannot afford a single listing has zero hidden blocking pairs and zero chance of placement under either policy. A group can consequently score well on a hidden-options measure while having almost no access to housing. We keep these households in the denominator, report their feasible sets, and examine how much discretion the platform has once tentative offers exist. Giving lower-income renters priority among equally large referral assignments had room to move only three additional placements toward them in the ordinary draws. A service that also changes affordability or expands search would face a different feasible graph.

\paragraph{Limits of the evidence.}
The housing data identify asking rents and surveyed household characteristics. Preferences, housing-provider priorities, budgets and assignments are modeled. The budget in the prompt is correlated with income, even though group membership is absent from the prompt. The language-model result conditions on one model version and one saved response per queried renter. The independent catalog check evaluates only the first-five benchmark. The data contain no broker relationships or commissions, and the study does not identify steering. In Chile the rank-informed model already has access to families' ordered applications. Its evidence concerns curation of known choices. The actual admissions system uses reserves and family contingencies beyond the pooled-priority replay. Changing what a platform displays could change what people search for, how they report, who applies, or whether they accept. If a shortlist is merely advice and families can freely apply off-list, the assumed restriction need not bind. Neither study estimates those responses, preference intensity, or a causal effect of deployment. The baseline school replay disagrees with the official course or outside option for 3.4--4.6\% of records, much larger than the net placement change (about 0.07--0.10 percentage points by year). Restricting the \emph{accounting}, after clearing each entire cohort, to applicants whose full-list replay agreed with the historical outcome still gives net losses of 356, 400, 399 and 259. This locates the result; replay agreement for one person cannot establish that all competitors would behave as modeled. The Chilean application years may include some of the same children. We treat them as finite cohorts and report exact counts. The shared first-five benchmark uses the same rule in both settings. The site-specific recommenders address different decisions. We do not pool their effect sizes.

\paragraph{What a platform can audit.}
Record the available options and capacities, the two complete sets of proposed lists, the reference offer where one exists, and the assignment induced by each set of lists. Those records distinguish a successful recommendation to its intended recipient, a lost place on a list that never changed, and a person without any feasible alternative. The last case points to a constraint on supply. The first two call for evaluating recommendation and allocation together.


\paragraph{Reproducibility.}
The accompanying code regenerates the two studies' aggregate tables and figures. In NYC we checked 2,880 published market-by-group endpoint rows against the regenerated menus. The legacy reply cache stores listing indices without a digest of the original prompt payload. Its consistency with this reconstruction can be checked, while the exact original payload cannot be authenticated independently. In Chile we recomputed the blind, rank-informed and full-list assignments in the original two cohorts and checked all 3,321,732 applicant--policy outcomes against the archived serial solver. The 2022--23 sign criterion was recorded locally before data acquisition. The extension to 2024--25 was committed to the research repository before downloading either release. The joint transitions, alternative-baseline, rollout, sensitivity and referral analyses are follow-ups. The two additional cohorts yielded a further 2,812,215 person--policy checks. All shared school outputs are aggregates.

\paragraph{Use of AI.} The team used OpenAI Codex extensively in developing the mathematical arguments, implementing and checking the replays, and drafting and revising the manuscript. The Claude model used to select rental menus is part of the experiment, a separate use described in Section~\ref{sec:twonyc}.
